% Composition study only — does not replace the reviewed Fig5 source.
% Narrative: preparation -> isolation -> reverse-force condition -> response.
\documentclass[border=4pt]{standalone}
\usepackage{polymer-paper-preamble}
\usetikzlibrary{arrows.meta,calc,positioning}

\begin{document}
\resizebox{178mm}{!}{%
\begin{tikzpicture}[
  every node/.style={font=\sffamily\fontsize{5.4}{6.5}\selectfont},
  panelLetter/.style={font=\sffamily\bfseries\fontsize{8}{9.6}\selectfont,text=cGray!92!black},
  panelTitle/.style={font=\sffamily\bfseries\fontsize{6.5}{7.8}\selectfont,text=cGray!86!black},
  label/.style={font=\sffamily\fontsize{5.2}{6.2}\selectfont,text=cGray!85!black},
  mute/.style={font=\sffamily\itshape\fontsize{5.1}{6.1}\selectfont,text=cGray!62!black},
  apparatus/.style={draw=cGray!76!black,fill=cGray!18,line width=.75pt},
  wire/.style={draw=cGray!78!black,line width=.75pt,line cap=round},
  beamOuter/.style={cAmber!74!black,line width=2.20mm,line cap=round},
  beamInner/.style={cAmber!10,line width=1.42mm,line cap=round},
  leader/.style={cGray!58!black,line width=.38pt,line cap=round},
  stageArrow/.style={-{Stealth[length=3.6pt,width=2.6pt]},cGray!62!black,line width=.62pt},
  force/.style={-{Stealth[length=4.4pt,width=3.2pt]},cRed!86!black,line width=.96pt},
  baseline/.style={-{Stealth[length=3.4pt,width=2.4pt]},cGray!62!black,line width=.58pt},
  q/.style={circle,fill=cRed!80!black,draw=cRed!95!black,line width=.26pt,minimum size=1.15mm,inner sep=0pt},
  terminal/.style={circle,draw=cGray!76!black,fill=white,line width=.52pt,minimum size=2.7pt,inner sep=0pt},
  axis/.style={-{Stealth[length=3.2pt,width=2.3pt]},cGray!72!black,line width=.60pt},
]

% A — actuation and charge formation.
\begin{scope}[shift={(0,4.55)}]
  \node[panelLetter,anchor=west] at (.10,3.17) {a};
  \node[panelTitle,anchor=west] at (.45,3.17) {actuation charge};
  \node[mute,anchor=west] at (.45,2.83) {field-on actuation};
  \fill[apparatus] (1.18,2.17) rectangle (2.28,2.48);
  \coordinate (a-root) at (1.73,2.17);
  \draw[wire] (a-root)--(1.73,2.72);
  \node[mute,anchor=west] at (2.00,2.69) {clip: GND};
  \draw[beamOuter] (a-root) .. controls (1.73,1.70) and (2.08,.78) .. (2.82,.42);
  \draw[beamInner] (a-root) .. controls (1.73,1.70) and (2.08,.78) .. (2.82,.42);
  \fill[apparatus] (4.40,.22) rectangle (4.72,2.48);
  \draw[leader] (4.56,2.48)--(4.56,2.76);
  \node[label,text=cRed!82!black,anchor=south] at (4.56,2.76) {$+5\,\mathrm{kV}$};
  \node[q] at (1.78,1.64) {}; \node[q] at (2.00,1.10) {}; \node[q] at (2.42,.61) {};
  \node[label,text=cRed!80!black,anchor=east] at (1.18,1.26) {trapped $q_{\mathrm{tr}}$};
  \draw[leader,cRed!58!black] (1.24,1.30)--(1.78,1.64);
  \draw[force] (1.87,1.46)--(3.58,1.46);
  \node[label,text=cRed!84!black,anchor=south] at (2.72,1.51) {attraction};
  \draw[cGray!50!black,line width=.38pt] (3.14,.33)--(3.14,.56);
  \draw[cGray!50!black,line width=.38pt] (4.32,.33)--(4.32,.56);
  \draw[{Stealth}-{Stealth},cGray!58!black,line width=.55pt] (3.18,.44)--(4.28,.44);
  \node[mute,anchor=south] at (3.73,.54) {air gap};
\end{scope}

% B — isolation as a full state, not an arrow caption.
\begin{scope}[shift={(6.05,4.55)}]
  \node[panelLetter,anchor=west] at (.10,3.17) {b};
  \node[panelTitle,anchor=west] at (.45,3.17) {source-off isolation};
  \node[mute,anchor=west] at (.45,2.83) {manual lead lift; clip floating};
  \fill[apparatus] (1.18,2.17) rectangle (2.28,2.48);
  \coordinate (b-root) at (1.73,2.17);
  \draw[wire] (b-root)--(1.73,2.61);
  \node[mute,anchor=west] at (2.00,2.59) {clip mounted};
  \draw[beamOuter] (b-root) .. controls (1.73,1.66) and (1.86,.74) .. (2.32,.28);
  \draw[beamInner] (b-root) .. controls (1.73,1.66) and (1.86,.74) .. (2.32,.28);
  \fill[apparatus] (4.40,.22) rectangle (4.72,2.48);
  \node[q] at (1.78,1.54) {}; \node[q] at (1.92,.96) {}; \node[q] at (2.16,.46) {};
  \draw[force,opacity=.60] (1.89,1.39)--(3.49,1.39);
  \node[label,text=cRed!72!black,anchor=south] at (2.69,1.44) {residual attraction};
  \draw[wire] (.78,.52)--(1.23,.52); \node[terminal] at (1.23,.52) {};
  \draw[wire] (1.83,.80)--(2.28,.80); \node[terminal] at (1.83,.80) {};
  \draw[baseline] (1.53,.57)--(1.53,.74);
  \node[mute,anchor=south] at (1.53,.91) {lead lifted};
  \node[label,anchor=west] at (.78,.18) {source OFF};
  \node[mute,anchor=west] at (2.55,.18) {electrically floating};
\end{scope}

% C — conditional force competition, made the last stage before the trace.
\begin{scope}[shift={(12.10,4.55)}]
  \node[panelLetter,anchor=west] at (.10,3.17) {c};
  \node[panelTitle,anchor=west] at (.45,3.17) {reversed drive};
  \node[mute,anchor=west] at (.45,2.83) {conditional force balance};
  \fill[apparatus] (1.18,2.17) rectangle (2.28,2.48);
  \coordinate (c-root) at (1.73,2.17);
  \draw[wire] (c-root)--(1.73,2.61);
  \node[mute,anchor=west] at (2.00,2.59) {floating clip};
  \draw[beamOuter] (c-root) .. controls (1.73,1.66) and (1.38,.74) .. (.80,.28);
  \draw[beamInner] (c-root) .. controls (1.73,1.66) and (1.38,.74) .. (.80,.28);
  \fill[apparatus] (4.40,.22) rectangle (4.72,2.48);
  \draw[leader] (4.56,2.48)--(4.56,2.76);
  \node[label,text=cRed!82!black,anchor=south] at (4.56,2.76) {$-5\,\mathrm{kV}$};
  \node[q] at (1.62,1.44) {}; \node[q] at (1.30,.84) {}; \node[q] at (.88,.34) {};
  \node[label,text=cRed!80!black,anchor=east] at (.98,1.62) {$q_{\mathrm{tr}}$};
  \draw[leader,cRed!58!black] (1.04,1.57)--(1.62,1.44);
  \draw[baseline] (1.74,1.52)--(2.74,1.52);
  \node[mute,anchor=south west] at (2.80,1.56) {Maxwell};
  \draw[force] (1.30,.84)--(.18,.84);
  \node[label,text=cRed!84!black,anchor=south east] at (.95,.89) {Coulomb};
  \node[label,anchor=north] at (2.70,.05) {reverse if $|F_{\mathrm C}|>|F_{\mathrm M}|$};
\end{scope}

% Reading-order connectors; they bind stages without introducing a separate mechanism.
\draw[stageArrow] (5.44,6.05)--(5.86,6.05);
\draw[stageArrow] (11.49,6.05)--(11.91,6.05);
\draw[stageArrow] (14.70,4.36)--(14.70,3.76);
\node[mute,anchor=west] at (14.82,4.08) {response};

% D — the qualitative response is wide enough to be read as the outcome of all stages.
\begin{scope}
  \node[panelLetter,anchor=west] at (.10,3.20) {d};
  \node[panelTitle,anchor=west] at (.45,3.20) {continuous bend response};
  \node[mute,anchor=west] at (.45,2.87) {qualitative trace from actuation onset};
  \draw[axis] (1.25,.60)--(1.25,2.55);
  \draw[axis] (1.25,1.52)--(17.10,1.52);
  \node[mute,rotate=90,anchor=south] at (.82,1.58) {bend angle};
  \node[mute,anchor=north east] at (17.10,.48) {time};
  \node[mute,anchor=east] at (1.17,2.30) {$+\theta$};
  \node[mute,anchor=east] at (1.17,.83) {$-\theta$};
  \draw[cBlue!82!black,line width=1.04pt]
    (1.38,1.52) .. controls (1.62,1.64) and (1.98,2.20) .. (2.38,2.30)
    -- (6.55,2.30)
    .. controls (6.88,2.30) and (7.02,2.18) .. (7.18,1.52)
    .. controls (7.31,.94) and (7.47,.67) .. (8.10,.61)
    .. controls (8.75,.56) and (10.25,.86) .. (11.40,1.18)
    .. controls (12.70,1.48) and (14.12,1.50) .. (15.86,1.50);
  \draw[cBlue!55!white,dashed,line width=.80pt]
    (15.86,1.50) .. controls (16.34,1.50) and (16.75,1.51) .. (17.10,1.52);
  \draw[cGray!45!black,dashed,line width=.50pt] (6.55,.62)--(6.55,2.30);
  \fill[cRed!85!black] (6.55,2.30) circle (.73pt);
  \draw[cGray!70!black,line width=.45pt] (1.38,1.52)--(1.38,1.34);
  \node[mute,anchor=north] at (1.38,1.28) {$t=0$};
  \node[label,anchor=south] at (4.38,2.38) {sustained attraction};
  \node[label,anchor=south] at (6.55,2.72) {source OFF};
  \node[mute,anchor=north] at (6.55,2.52) {clip floating};
  \node[label,text=cRed!84!black,anchor=south west] at (7.32,1.77) {reversed drive};
  \node[mute,anchor=south west] at (11.60,1.22) {slow recovery};
\end{scope}
\end{tikzpicture}%
}
\end{document}
